\section{Strict retractions and prescribed fence homotopies}
\label{sec:D}

All categories in this section are $\mathcal U$-small. A categorical homotopy is a finite zigzag of natural transformations. We write $Z$ for the free category on a finite oriented interval, and distinguish its initial vertex $0$ and final vertex $z$ (these names specify positions, not initial or terminal objects in the category). A padding of $Z$ is a longer fence $Z'$ with a functor $Z'\to Z$ preserving these endpoint positions; added edges may map to identities.

\begin{claim}{D.1}{Proved}{Strict retraction in the prescribed weak-homotopy class}
Let $A$ be Thomason-fibrant, let $i:A\hookrightarrow B$ be a Thomason cofibration, and let $p:B\to A$ satisfy $pi\simeq\mathrm{id}_A$ by a finite categorical homotopy. There is $r:B\to A$ with $ri=\mathrm{id}_A$ in $\mathrm{Cat}$ and $[r]=[p]$ in $\mathrm{Hot}$.
\end{claim}
\begin{proof}
Represent the given homotopy by $H:A\times Z\to A$, with $H(-,0)=pi$ and $H(-,z)=\mathrm{id}_A$. The constant-path functor $c:A\to\mathrm{Fun}(Z,A)$ is a categorical homotopy equivalence: delete the endpoints of the fence successively, using each incident arrow for the natural transformation between the identity and the resulting endpoint-collapse endofunctor, and compose the resulting finite zigzags. Consequently $c$ is a Thomason weak equivalence.

Factor the endpoint evaluation functor as
\[
 \mathrm{Fun}(Z,A)\xrightarrow{j}P
 \xrightarrow{(d_0,d_1)}A\times A,
\]
with $j$ a trivial cofibration and the second map a fibration. The maps $d_0,d_1:P\to A$ are fibrations because the projections $A\times A\to A$ are pullbacks of the fibration $A\to *$. The composite $jc$ is a weak equivalence and $d_0jc=d_1jc=\mathrm{id}_A$, so both $d_0,d_1$ are weak equivalences as well. Thus both are trivial fibrations. Let $h:A\to P$ be $j$ applied to the adjoint of $H$. The square with top $h$, left $i$, bottom $p$, and right $d_0$ commutes. Lift it to $\widetilde h:B\to P$ and set $r=d_1\widetilde h$. The equations $ri=d_1h=\mathrm{id}_A$ and $d_0\widetilde h=p$ hold in $\mathrm{Cat}$. In $\mathrm{Hot}$ both $d_0$ and $d_1$ equal $(jc)^{-1}$; hence $[r]=[p]$ there.
\end{proof}

This confirms every step of the weaker statement in CATHOT E.6. The factorization supplies no finite natural-transformation zigzag on $B$ between $r$ and $p$.

\begin{claim}{D.2}{Proved}{The local-target and groupoid cases}
Under the hypotheses of D.1, a local target $A$ admits $r$ with $[r]=[p]$ in $\mathrm{Cat}_{\mathrm{hot}}$. For an actual groupoid target the same conclusion follows for every full inclusion $i$ that is injective on objects, without Thomason fibrancy or cofibration assumptions.
\end{claim}
\begin{proof}
For the first assertion, locality makes the comparison
$\pi_0\mathrm{Fun}(B,A)\to[BB,BA]$ injective, and D.1 supplies equal images.
For the second, the given zigzag composes, after reversing its invertible components where necessary, to a natural isomorphism $\alpha:pi\Rightarrow\mathrm{id}_A$. Identify $A$ with its full image in $B$. For $b\in A$ put $r(b)=b$ and $\beta_b=\alpha_b:p(b)\to b$; for $b\notin A$ put $r(b)=p(b)$ and $\beta_b=\mathrm{id}$. For $u:b\to b'$ define
$r(u)=\beta_{b'}p(u)\beta_b^{-1}$.
The identities and composition equations follow by cancellation of consecutive $\beta_b^{-1}\beta_b$. Naturality of $\alpha$ proves $ri=\mathrm{id}_A$ on all arrows as well as objects. The family $\beta$ is a natural isomorphism $p\Rightarrow r$.
\end{proof}

\begin{claim}{D.3}{Proved}{Prescribed fence extension for targets with a terminal object}
Let $i:C\hookrightarrow D$ be a full sieve inclusion, let $F$ have a terminal object $t$, and let $f:D\to F$ and $H:C\times Z\to F$ agree on $C\times\{0\}$. There is an extension over $D\times Z'$ for a padding $Z'\to Z$ obtained by adding one edge at the initial position.
\end{claim}
\begin{proof}
Define $f':D\to F$ by $f'|_C=f|_C$ and $f'(d)=t$ outside $C$. Since $C$ is a sieve, no arrow has source outside $C$ and target in $C$. Every remaining mixed arrow maps to the unique arrow into $t$. The identity components on $C$ and the terminal arrows outside $C$ define $f\Rightarrow f'$. Next define $\widehat H:D\times Z\to F$ by $H$ on $C\times Z$ and the constant value $t$ outside. Mixed naturality squares commute by uniqueness of arrows into $t$. Prepend the transformation $f\Rightarrow f'$ and map the added edge to the identity at $0\in Z$. Its restriction to $C$ is stationary, so the resulting extension has exactly the required boundary.
\end{proof}

\begin{claim}{D.4}{Proved}{Prescribed fence extension for targets with an initial object}
Let $i:C\hookrightarrow D$ be a pseudo-Dwyer inclusion. Choose its cosieve neighborhood $W\subseteq D$, a retraction $s:W\to C$, and a transformation $\varepsilon:is\Rightarrow\mathrm{id}_W$ whose restriction to $C$ is the identity. Let $F$ have an initial object $o$. Every pair $f:D\to F$, $H:C\times Z\to F$ agreeing at the initial endpoint admits the extension of D.3 after one padding edge. When $W=D$, the conclusion holds for every target $F$.
\end{claim}
\begin{proof}
On $W$ put $f'=fis$, and outside $W$ put $f'=o$. There are no arrows from $W$ to its complement because $W$ is a cosieve. Define arrows with source outside $W$ and target in $W$ using the unique maps from $o$. This specifies a functor, and $f\varepsilon$ on $W$, together with the initial arrows elsewhere, specifies $f'\Rightarrow f$. Its restriction to $C$ is the identity. Extend $H$ to $D\times Z$ by $H(s(-),-)$ on $W$ and the constant initial object on the complement. The same cosieve and uniqueness checks prove functoriality and naturality. Prepend the reversed-in-position edge $f\Leftarrow f'$, mapped to the identity at $0\in Z$. This is the desired padding. With $W=D$ the construction has no complementary part and uses no initial object.
\end{proof}

\begin{claim}{D.5}{Proved}{A non-cofibrant corner can still admit padded extension}
Let $C=\{0\}$, $D=[1]$, and let $F$ be the three-object poset with arrows $o\to u$ and $o\to v$. The corner functor sending $(0,0)$ to $o$, $(1,0)$ to $u$, and $(0,1)$ to $v$ has no extension to $D\times[1]$. It extends over a padded fence with two edges. The category $F$ is Thomason-fibrant.
\end{claim}
\begin{proof}
An extension over the square requires an upper bound of $u,v$, which does not exist in $F$. Over the fence $0\leftarrow0'\to1$, mapping its first edge to $\mathrm{id}_0$ in $[1]$, take the successive functors $D\to F$ to be $(o\to u)$, $(o\xrightarrow{\mathrm{id}}o)$, and $(v\xrightarrow{\mathrm{id}}v)$. The two transformations are $f_0\Leftarrow f_{0'}\Rightarrow f_1$, and restriction to $C$ is $o\Leftarrow o\Rightarrow v$, the required pullback of the original boundary.
To check fibrancy, apply the construction underlying D.4 to each generating Thomason horn inclusion: it is Dwyer, so its cosieve neighborhood has a right-adjoint retraction. Given a map from the horn to $F$, define its extension by that retraction on the neighborhood and by $o$ elsewhere. The same arrow check proves a strict extension. Thus $F\to *$ has the required right lifting property.
\end{proof}


\begin{claim}{D.6}{Refuted}{Fibrancy does not make the categorical comparison injective}
There are a finite category $X$ with six objects and twenty-two arrows, a $\mathcal U$-small Thomason-fibrant poset $F$, and two functors $u,v:X\to F$ such that $[u]\ne[v]$ in $\mathrm{Cat}_{\mathrm{hot}}$, while $Bu$ and $Bv$ are unpointedly homotopic. In particular fibrancy of the target does not imply injectivity of
\[
 \pi_0\mathrm{Fun}(X,F)\longrightarrow[BX,BF].
\]
\end{claim}
\begin{proof}
Let $P$ have two elements in each of three levels $0,1,2$, with $x<y$ exactly when the level of $x$ is smaller than the level of $y$. Its order complex is the simplicial join of three copies of $S^0$, so $BP$ is homeomorphic to $S^2$. Let $X$ be the free category on the Hasse quiver of $P$: there are four arrows from level $0$ to level $1$, four from level $1$ to level $2$, and eight length-two composites. Together with the six identities these are all twenty-two arrows. Let $q:X\to P$ send a path to its endpoints.

For every poset $T$, precomposition gives an isomorphism of categories
\[
 q^*:\mathrm{Fun}(P,T)\xrightarrow{\cong}\mathrm{Fun}(X,T).
\]
Indeed a functor on the free quiver specifies an order-preserving assignment on the eight generating comparisons. Transitivity forces exactly the four additional level-$0$ to level-$2$ comparisons; uniqueness of arrows in $T$ identifies the two composites between each such pair. A natural transformation on either side is exactly a pointwise comparison of the same object assignments, and the naturality squares commute because $T$ is a poset. This proves the claimed isomorphism on both objects and arrows.

Take a fibrant replacement $j:P\to F$ in the Thomason model structure on posets, whose existence is Raptis, Theorem~2.6. The inclusion of posets into categories creates fibrations and weak equivalences, so $F$ is Thomason-fibrant as a category and $Bj$ is a weak equivalence. Set $u=jq$ and let $v$ be any constant functor. An equality $[u]=[v]$ in $\mathrm{Cat}_{\mathrm{hot}}$ would, by the displayed isomorphism, give a finite natural-transformation zigzag from $j:P\to F$ to a constant functor. Applying the nerve and realization gives a homotopy between $Bj$ and a constant, contradicting the isomorphism $H_2(Bj;\mathbb Z):\mathbb Z\to\mathbb Z$.

On the other hand, $NX$ collapses onto the Hasse graph. Each of its eight nondegenerate two-simplices has its own composite long edge; that edge belongs to no other nondegenerate two-simplex. Delete each triangle together with its free long edge. The resulting connected graph has six vertices and eight edges, hence is a wedge of three circles. The space $BF$ is simply connected because it is weakly equivalent to $S^2$. Therefore every map from this graph to $BF$ is nullhomotopic: contract a spanning tree, and nullhomotope each of its three loop maps using $\pi_1BF=0$. It follows that $Bu\simeq Bv$ as claimed. All category compositions, free-face incidences, and the chain-complex ranks are checked in \texttt{code/d\_finite\_checks.py}; the infinite fibrant replacement is justified by the specified model structure, not by a bounded computation.
\end{proof}

\paragraph{Unsettled statement.}
D.1--D.6 do not establish the general Cat$_{\mathrm{hot}}$ assertion of Part D or the general prescribed-fence extension property for an arbitrary Thomason-fibrant target. No counterexample satisfying all those hypotheses has been obtained. The current exact obstruction is the production of a single finite fence on the entire source from the model-path lift in D.1; ordinary product-corner non-cofibrancy does not settle this, as D.5 shows. D.6 blocks an attempted proof by asserting homotopy-injectivity for all fibrant targets.
